%% sections/04-study-design.tex

\section{Study Design}
\label{sec:study}

\subsection{The Marathon Workshop Pipeline}

The Marathon D\&D Workshop is a structured pipeline for designing and playtesting
D\&D 5e adventures. The pipeline has seven stages:

\begin{enumerate}
  \item \textbf{PREP}: Adventure compiled to a module document; party character sheets
        loaded; rubric version recorded in session frontmatter.
  \item \textbf{PLAY}: AI system (Claude, Anthropic) runs the session as DM while
        simultaneously playing all PCs per their Playstyle Heuristics.
  \item \textbf{LOG}: Session log produced as a structured narrative document recording
        all scenes, dice rolls (with seed for reproducibility), PC actions, and
        narrative outcomes.
  \item \textbf{GATE}: Session log scored against the locked rubric version using the
        session-gate scoring protocol. Produces per-dimension scores and overall gate score.
  \item \textbf{PANEL}: Five-lens player-experience panel reviews the session log from
        five distinct player perspectives (Tactician, Roleplayer, Lorekeeper, Improviser,
        Fresh Face), producing weighted aggregate panel scores.
  \item \textbf{INNOVATE}: Session log and panel reviews scanned for innovations not
        captured by existing rubric anchors. Logged entries may trigger amendment proposals.
  \item \textbf{HANDOFF}: Canonical handoff document produced; TRACKER.md updated with
        session metrics; campaign-permanent facts registered.
\end{enumerate}

All dice rolls use a seeded pseudorandom number generator for reproducibility.
The seed is recorded in the session log header. This ensures that sessions can be
re-run from the same seed to confirm outcomes.

\subsection{The Seven Campaigns}

The corpus consists of 49 sessions across 7 complete campaigns, each comprising
7 adventures. Table~\ref{tab:campaigns} summarizes the campaigns.

\begin{table*}[t]
\caption{The Seven Marathon Campaigns: Party, Archetype, Central Question, Route, Mean Scores}
\label{tab:campaigns}
\begin{tabular}{@{}llllrrl@{}}
\toprule
Campaign & Party & Archetype & Central Question & Gate & Panel & Route \\
\midrule
C1 Moon-Silver Cycle & varduin-muster & grief-themed & What is an emotion worth trading? & 74.1 & 72.2 & D \\
C2 Conclave Compact & compact-wardens & covenant-restoration & When a covenant breaks, who mends it? & 76.3 & 74.2 & D \\
C3 Halted Spire & the-witnesses & faith/authority & Who decides who is worth saving? & 71.7 & 71.8 & D \\
C4 Thorngate Watch & the-relief & siege/resource & How much of yourself do you spend? & 74.0 & 74.4 & E \\
C5 Thieves Guild & the-guild & professional-code & When the guild becomes what you steal from? & 73.9 & 73.8 & B \\
C6 Military Unit & the-unit & duty-spine & When orders become wrong, how do you know? & 74.6 & 74.5 & E \\
C7 Party of Strangers & the-strangers & no-designed-stakes & What do people with nothing in common share? & 73.4 & 73.4 & E \\
\bottomrule
\end{tabular}
\end{table*}

Campaigns differ on several design axes. C1--C3 used grief-adjacent party archetypes
with explicit loss-hooks on every PC character sheet. C4 introduced resource exhaustion
(Supply Die + between-session HP carry-over; no full long rest for 7 sessions). C5
used professional-code motivation with no grief hooks. C6 used a duty-spine with an
institutional deliverable (a field journal presented verbatim at a tribunal). C7 used
parties of strangers with no designed personal stakes, no grief paragraphs, no loyalty
objects, and no professional codes.

\subsection{Session Log Format}

Each session log records: a PREP frontmatter block (adventure slug, party slug, rubric
version, dice seed); a sequential scene log (scene number, scene summary, all dice
rolls with natural values and modifiers, PC actions and narration); a SURPRISE log
(moments tagged as exceeding design prediction); and a POST-PLAY note (observed
innovations, amendment candidates, campaign-permanent fact registrations).

Gate scoring reads the session log and populates a scorecard with per-dimension
evidence cells (specific log citations) and scores. Advisory scores (sessions 1--3 per
campaign) are not treated as binding for pass/fail determination.

\subsection{Session-Gate Scoring Protocol}

For each dimension, the reviewer reads the session log with that dimension as focus,
finds supporting evidence (quoted passages or scene references), finds counter-evidence,
assigns a score using the rubric's band descriptions, and writes an evidence cell.
The scoring protocol explicitly prohibits generous or harsh adjustment: ``Do not score
generously out of sympathy. Do not score harshly to seem rigorous. Score what is in
the log.''

Gate scores and panel scores are produced independently. Panel scores use the five-lens
weighted aggregate (each lens scores 0--10 on its own priority axes; the aggregate is
weighted by the lens's contribution to the campaign archetype's design). The gate/panel
spread (difference between gate and panel scores) is monitored as a diagnostic: large
spreads indicate a structural artifact (e.g., a session with zero combat dice produces a
panel dip from the Tactician lens because the lens has fewer events to evaluate).
